% --- Archon physics formalization source begin ---
% archon:physics
% archon:chemistry
% archon:covers IChO2026Problems/problem_icho_2026_t9_a6.lean
% archon:source-report reports/icho_2026/problem_icho_2026_t9_a6.source.json
% archon:problem-id icho_2026_t9
% archon:part-id T9-A6
% archon:previous-part icho_2026_t9_a5
% archon:previous-part-policy natural_language_prerequisite_only

\chapter{Icho Chemistry problem icho\_2026\_t9\_a6}
\label{ch:IChO2026Problems_problem_icho_2026_t9_a6}

\paragraph{Problem source.}
\#\# Source
58th International Chemistry Olympiad, Tashkent, Uzbekistan, 2026.
Official-materials index: https://www.icho2026.uz/problems
English problem PDF: ../icho\_2026\_source/raw/theory\_problem.pdf
English solution/rubric PDF: ../icho\_2026\_source/raw/theory\_solution.pdf
Problem PDF page 86; printed local question page 3; solution starts on PDF page 86.
The page PNG is primary visual evidence for formulas, structures, tables, and figures.

\#\# T9. Cyclodextrin Chemistry
T9. Cyclodextrin Chemistry 7\% Cyclodextrins (CD) are a family of cyclic oligosaccharides, consisting of glucose subunits joined by α1,4-glycosidic bonds. The three most common cyclodextrins (α-, β-, γ-cyclodextrin) contain 6, 7, and 8 α-D-glucopyranoside units, respectively. 9.1 Calculate the molar mass of β-CD. Assume M𝑤(glucose) = 180.16 g mol−1. CDs combine macrocyclic properties with glucose chemistry, which was demonstrated in the synthesis of X. Primary and secondary hydroxy groups show different reactivity due to being a part of CDs. Stoddart et al. synthesised compound K, where only one OH group remained free per glucopyranoside unit. equiv.= equivalent 9.2 Tick the favorable chair conformation and draw the structure of K by completing the CD template. Show stereochemistry unambiguously. 9.3 Determine the ring size, 𝑟𝑠, of macrocycle X. Give the number of stereocentres, 𝑠𝑐, in X. α-Cycloaltrin, a synthetic analogue of α-CD, was synthesised as follows. 9.4 Draw the structure of Y with stereochemistry by completing the CD template. In 2000, Sinay et al. demonstrated that a single protic group (NH and OH) at unit 1 directs the next reductive debenzylation of a primary OH group to unit 4 in the macrocyclic ring (or to unit 3 if the unit 4 position is not available). This allowed the synthesis of the following β-CD dimer as a mixture of constitutional linkage isomers. 9.5 Draw the structure of L on the β-CD template. Fill in all the boxes.

\#\# Current subquestion 9.6
Determine the number of isomers of the β-CD dimer that can form during the synthesis by Sinay et. al.

\#\# Dataset metadata
IChO 2026; T9; raw rubric points=4; kind=theory; formalization\_ready=true.

\paragraph{Shared problem context.}
T9. Cyclodextrin Chemistry 7\% Cyclodextrins (CD) are a family of cyclic oligosaccharides, consisting of glucose subunits joined by α1,4-glycosidic bonds. The three most common cyclodextrins (α-, β-, γ-cyclodextrin) contain 6, 7, and 8 α-D-glucopyranoside units, respectively. 9.1 Calculate the molar mass of β-CD. Assume M𝑤(glucose) = 180.16 g mol−1. CDs combine macrocyclic properties with glucose chemistry, which was demonstrated in the synthesis of X. Primary and secondary hydroxy groups show different reactivity due to being a part of CDs. Stoddart et al. synthesised compound K, where only one OH group remained free per glucopyranoside unit. equiv.= equivalent 9.2 Tick the favorable chair conformation and draw the structure of K by completing the CD template. Show stereochemistry unambiguously. 9.3 Determine the ring size, 𝑟𝑠, of macrocycle X. Give the number of stereocentres, 𝑠𝑐, in X. α-Cycloaltrin, a synthetic analogue of α-CD, was synthesised as follows. 9.4 Draw the structure of Y with stereochemistry by completing the CD template. In 2000, Sinay et al. demonstrated that a single protic group (NH and OH) at unit 1 directs the next reductive debenzylation of a primary OH group to unit 4 in the macrocyclic ring (or to unit 3 if the unit 4 position is not available). This allowed the synthesis of the following β-CD dimer as a mixture of constitutional linkage isomers. 9.5 Draw the structure of L on the β-CD template. Fill in all the boxes.

\paragraph{Current subquestion.}
Determine the number of isomers of the β-CD dimer that can form during the synthesis by Sinay et. al.

\paragraph{Recorded answer/context.}
Due to β-CD consisting of seven units and being chiral, two regioisomers are formed in the alkylation reaction (M-4R, M-1R). Therefore, the dimer would have three isomers with following connections: 4-4, 1-1, 1-4. (Note 4-1 = 1-4) Correct number of isomers 4 points. 4.0 pt

\paragraph{Figure/image paths.}
\begin{itemize}
\item ../icho\_2026\_source/image/T9\_page-3.png
\item ../icho\_2026\_source/image/T9\_page-2.png
\end{itemize}

\paragraph{Reusable previous-part conclusions.}
\begin{itemize}
\item Source T9-A5. Question: Draw the structure of L on the β-CD template. Fill in all the boxes. Reusable conclusions: Visual-answer notice: the official solution contains essential ticks, structures, tables, graphs, or equations that PDF plain-text extraction may omit. Consult the cited solution PDF page.

Correct structure 2 points. -1 p for each incorrect substituent 0 p if 2 or more incorrect substituents 2.0 pt Policy: natural\_language\_prerequisite\_only; the current target and later answers must not be assumed
\end{itemize}

\paragraph{Formalization target.}
The β-CD ring is modeled as \texttt{ZMod 7} with primary-face substitution
patterns valued in \texttt{\{OBn, OH, alkenyl\}}, identified up to rotation only
(chirality: no reflection).  The Sinay directing rule places the second free OH
of the intermediate L at ring offset $+3$ from the directing unit; alkylating
either free OH of L gives the two monomer regioisomers M-1R and M-4R, and a
dimer is an unordered pair (\texttt{Sym2}) of monomer regioisomers.

\begin{definition}[β-CD ring, Sinay offset, L, alkylated monomers, dimers]
\label{def:physics:icho_2026_t9_a6:setup}
\lean{IChO2026.T9.A6.betaCDNumUnits, IChO2026.T9.A6.UnitPos, IChO2026.T9.A6.PrimarySubstituent, IChO2026.T9.A6.Pattern, IChO2026.T9.A6.RotationEquiv, IChO2026.T9.A6.sinayOffset, IChO2026.T9.A6.IsLPattern, IChO2026.T9.A6.IsAlkylatedMonomer, IChO2026.T9.A6.MonomerIsomer, IChO2026.T9.A6.DimerIsomer, IChO2026.T9.A6.patternL, IChO2026.T9.A6.patternM1R, IChO2026.T9.A6.patternM4R, IChO2026.T9.A6.M1R, IChO2026.T9.A6.M4R}
\leanok
Patterns of primary substituents on the seven units of β-CD, rotation
equivalence, the Sinay offset ($+3$ for ring size $7$), the intermediate L
(two free primary OH groups at offset $+3$), single-site alkylated monomers,
monomer regioisomers as rotation classes, and dimers as unordered pairs of
monomer regioisomers.
\end{definition}

\begin{lemma}[The two monomer regioisomers]
\label{lem:physics:icho_2026_t9_a6:monomers}
\lean{IChO2026.T9.A6.isLPattern_patternL, IChO2026.T9.A6.isAlkylatedMonomer_patternM1R, IChO2026.T9.A6.isAlkylatedMonomer_patternM4R, IChO2026.T9.A6.M1R_ne_M4R, IChO2026.T9.A6.monomer_isomers_enumeration, IChO2026.T9.A6.monomer_regioisomers_card}
\leanok
The canonical L satisfies the L-specification; M-1R and M-4R arise by
alkylating L at unit 1 resp.\ unit 4; they are distinct because a rotation
identifying them would have to send the alkenyl site $0$ to $3$, i.e.\ be the
shift $+3$, which sends the free-OH site $3$ to $6 \neq 0$
(equivalently $+3 \not\equiv -3 \pmod 7$); every valid alkylated monomer is
M-1R or M-4R; hence exactly two monomer regioisomers exist.
\end{lemma}
\begin{proof}
\leanok
Validity of L and of the two alkylations is by finite case analysis on
\texttt{ZMod 7}.  For the enumeration, an L-pattern is determined up to
rotation by its OH pair $\{i, i+3\}$, and the alkylated site is one of the two
free OH positions; composing with the rotation $-i$ matches the monomer onto
\texttt{patternM1R} or \texttt{patternM4R}.  The cardinality is transported
along the bijection $\texttt{Bool} \to $ monomer isomers.
\end{proof}

\begin{theorem}[IChO 2026 T9.6: number of dimer isomers]
\label{thm:physics:icho_2026_t9_a6:target}
\lean{IChO2026.T9.A6.dimer_isomers_enumeration, IChO2026.T9.A6.dimer_isomers_distinct, IChO2026.T9.A6.icho_2026_t9_a6_number_of_dimer_isomers}
\leanok
The number of constitutional (linkage) isomers of the β-CD dimer that can form
during the synthesis by Sinay et al.\ is $3$: the 4-4, 1-1 and 1-4 connected
dimers (note $4\text{-}1 = 1\text{-}4$).
\end{theorem}
\begin{proof}
\leanok
Every dimer is an unordered pair of monomer regioisomers, hence one of
$s(\mathrm{M4R},\mathrm{M4R})$, $s(\mathrm{M1R},\mathrm{M1R})$,
$s(\mathrm{M1R},\mathrm{M4R})$ by the monomer enumeration and symmetry of
\texttt{Sym2}.  The three are pairwise distinct by the component
characterization of \texttt{Sym2} equality together with
$\mathrm{M1R} \neq \mathrm{M4R}$.  The count is transported along the
bijection $\texttt{Fin 3} \to$ dimer isomers.
\end{proof}
% --- Archon physics formalization source end ---
